% ============================================================================
%  SECTION 5 — ISTA : ITERATIVE SHRINKAGE-THRESHOLDING
%  >>> A REMPLIR PAR LE MEMBRE 3 <<<
%
%  Depend de la section 4 (operateur proximal) du Membre 2.
%
%  Contenu attendu (voir REPARTITION_DETAILLEE.md) :
%    - Construction d'ISTA comme MAJORATION-MINIMISATION : a chaque pas, on
%      minimise un majorant quadratique de f + le terme g ; montrer que cela
%      redonne exactement "pas de gradient + prox".
%    - Iteration :
%        x_{k+1} = \softthr( x_k - t \nabla f(x_k), t\lambda )
%      avec \nabla f(x) = A^T(Ax - y) et pas t <= 1/L, L = ||A^T A||.
%    - PREUVE DE CONVERGENCE EN O(1/k) (argument du gradient Lipschitz).
%
%  Penser a renvoyer a la figure de convergence de la section 9 (experiences).
% ============================================================================
\section{ISTA : l'algorithme de gradient proximal}
\label{sec:ista}

% --- [MEMBRE 3] Ecrire la section ici. Squelette suggere ci-dessous. --------

\subsection{Idee : separer la partie lisse et la partie non lisse}
% TODO (Membre 3)

\subsection{Construction par majoration-minimisation}
% TODO (Membre 3)

\subsection{L'algorithme}
% TODO (Membre 3) : pseudo-code + role du pas t <= 1/L.

\subsection{Convergence en $O(1/k)$}
% TODO (Membre 3) : enonce + preuve.
